%! Characteristics of Radical Functions — Unit 5, Lesson 5.0 — Group Activity (Answer Key)
\documentclass[10pt]{article}
\usepackage{algebra2-article}
\usepackage{algebra2-key}   % pulls in -boxes; do NOT also load -boxes
\newcommand{\TallMath}[1]{$\displaystyle #1\rule[-1.4em]{0pt}{3.2em}$}
\newcommand{\lgrid}[4]{%
  \draw[step=1, linegray!40, very thin] (#1,#3) grid (#2,#4);
  \draw[->, semithick, charcoal] (#1,0)--(#2,0) node[below right, font=\scriptsize]{$x$};
  \draw[->, semithick, charcoal] (0,#3)--(0,#4) node[left, font=\scriptsize]{$y$};}
\newcommand{\cbrtpos}[2]{\draw[very thick, forest, domain=#1:#2, samples=70, smooth]
  plot (\x,{pow(\x,1/3)-1});}
\newcommand{\cbrtneg}[2]{\draw[very thick, forest, domain=#1:#2, samples=70, smooth]
  plot (\x,{-pow(-1*(\x),1/3)-1});}

\begin{document}

\pageheader{Unit 5, Lesson 5.0}{Group Activity --- Answer Key}
\namepartnerperiod

\vspace{0.08in}

\begin{headlinebox}{forestbg}
\small\textbf{Reading Radical Graphs.} A radical function tells you \emph{where it is allowed to
exist} before it tells you anything else --- and the \textbf{index} is what decides. With your
partner, read each graph carefully and \emph{justify} your answers. Write every domain and range in
\textbf{interval notation}, and check each one against the radicand.
\end{headlinebox}

\vspace{0.06in}

\begin{tcolorbox}[colback=white, colframe=black!40, title=\textbf{Tier R --- Remediate},
    fonttitle=\bfseries, arc=1.5mm, left=3mm, right=3mm, top=2mm, bottom=2mm, breakable]
The graph shows $g(x) = \sqrt{x-1}$.
\begin{center}
\begin{tikzpicture}[scale=0.46]
  \lgrid{-3}{11}{-2}{4}
  \begin{scope}
    \clip (-3,-2) rectangle (11,4);
    \draw[very thick, forest, domain=1:11, samples=100, smooth] plot (\x,{sqrt((\x)-1)});
  \end{scope}
  \fill[forest] (1,0) circle (4pt) node[below left, font=\scriptsize]{$(1,0)$};
  \fill[charcoal] (2,1) circle (2.6pt);
  \fill[charcoal] (5,2) circle (2.6pt) node[above left, font=\scriptsize]{$(5,2)$};
  \fill[charcoal] (10,3) circle (2.6pt);
\end{tikzpicture}
\end{center}
\begin{enumerate}[leftmargin=1.7em, itemsep=7pt]
  \item The \textbf{radicand} is \ans{$x-1$}. Set it $\ge 0$: \ans{$x\ge 1$}, so the
        \textbf{domain} is \ans{$[1,\infty)$}.
  \item The graph begins at a single point called the \textbf{endpoint}: \ans{$(1,0)$}. \\
        It is drawn as a \emph{solid} dot. Why is that input included in the domain?
        \ans{Because $x=1$ makes the radicand $0$, and $\sqrt{0}=0$ is a perfectly good real number.}
  \item \textbf{Range:} \ans{$[0,\infty)$} \quad \textbf{$x$-intercept:} \ans{$(1,0)$}
  \item Try to find the \textbf{$y$-intercept} by substituting $x=0$. The radicand becomes
        \ans{$-1$}, so $g(0)$ is \ans{not a real number}. \\
        What does that tell you about whether this graph has a $y$-intercept?
        \ans{$\sqrt{-1}$ does not exist, so $x=0$ is outside the domain. The graph never reaches
        the $y$-axis, so there is \emph{no} $y$-intercept.}
  \item $g$ is \textbf{increasing} on \ans{$[1,\infty)$}. \quad It has an absolute
        \ans{minimum} of \ans{0}, and \emph{no} absolute \ans{maximum}.
  \item \textbf{End behavior:} as $x\to+\infty$, $g(x)\to$ \ans{$+\infty$}. \\
        As $x\to-\infty$, $g(x)\to$ \dots\ careful! Explain what actually happens on the left.
        \ans{Nothing --- there is no left end. The graph stops at $x=1$ because everything to the
        left of it is outside the domain.}
\end{enumerate}
\end{tcolorbox}

\begin{tcolorbox}[colback=white, colframe=black!40, title=\textbf{Tier A --- Approaching Proficiency},
    fonttitle=\bfseries, arc=1.5mm, left=3mm, right=3mm, top=2mm, bottom=2mm, breakable]
Two radical functions --- \emph{and only one of them has a restricted domain}. Fill in the whole
table, then answer the two justification questions.
\begin{center}
\begin{tikzpicture}[scale=0.42]
  % p(x) = 2 - sqrt(x)
  \begin{scope}
    \lgrid{-2}{10}{-3}{3}
    \begin{scope}
      \clip (-2,-3) rectangle (10,3);
      \draw[very thick, forest, domain=0:10, samples=100, smooth] plot (\x,{2-sqrt(\x)});
    \end{scope}
    \fill[forest] (0,2) circle (4pt);
    \fill[charcoal] (1,1) circle (3pt);
    \fill[charcoal] (4,0) circle (3pt);
    \fill[charcoal] (9,-1) circle (3pt);
    \node[charcoal, font=\scriptsize] at (4,-4.4) {$p(x)=2-\sqrt{x}$};
  \end{scope}
  % q(x) = cbrt(x) - 1
  \begin{scope}[xshift=15cm]
    \lgrid{-9}{9}{-4}{2}
    \begin{scope}
      \clip (-9,-4) rectangle (9,2);
      \cbrtpos{0}{9}
      \cbrtneg{-9}{0}
    \end{scope}
    \fill[charcoal] (-8,-3) circle (3pt);
    \fill[charcoal] (0,-1) circle (3pt);
    \fill[charcoal] (1,0) circle (3pt);
    \fill[charcoal] (8,1) circle (3pt);
    \node[charcoal, font=\scriptsize] at (0,-5.4) {$q(x)=\sqrt[3]{x}-1$};
  \end{scope}
\end{tikzpicture}
\end{center}
\vspace{0.20in}
\renewcommand{\arraystretch}{1.5}
\begin{tabularx}{\linewidth}{@{}>{\bfseries}p{0.32\linewidth} X X@{}}
 & \textbf{$p(x)=2-\sqrt{x}$} & \textbf{$q(x)=\sqrt[3]{x}-1$} \\
\hline
Index (even or odd) & \ans{even ($2$)} & \ans{odd ($3$)} \\
Domain & \ans{$[0,\infty)$} & \ans{$(-\infty,\infty)$} \\
Range & \ans{$(-\infty,2]$} & \ans{$(-\infty,\infty)$} \\
Endpoint (if any) & \ans{$(0,2)$} & \ans{none} \\
$x$-intercept & \ans{$(4,0)$} & \ans{$(1,0)$} \\
$y$-intercept & \ans{$(0,2)$} & \ans{$(0,-1)$} \\
Increasing or decreasing? & \ans{decreasing} & \ans{increasing} \\
Absolute max or min? & \ans{abs.\ max of $2$} & \ans{neither} \\
\end{tabularx}
\vspace{4pt}

\textbf{Justify 1.} $p$ has a restricted domain and $q$ does not. Point to the exact feature of each
\emph{equation} that causes this --- and explain why that feature has the effect it does.
\ansline{$p$ has an \emph{even} index (the plain $\sqrt{\phantom{x}}$), so its radicand $x$ must be
$\ge 0$ --- no real number squares to a negative --- giving the domain $[0,\infty)$. $q$ has an
\emph{odd} index ($\sqrt[3]{\phantom{x}}$), and cubing preserves sign, so every real number is a
legal radicand and the domain is all reals. It is the \textbf{index}, not the radical symbol, that
restricts.}

\textbf{Justify 2.} $p$ has an absolute \emph{maximum}, but the parent $y=\sqrt{x}$ has an absolute
\emph{minimum}, and $q$ has neither. Explain all three at once, using the words \textbf{endpoint} and
\textbf{direction}. \ansline{A square root graph has an \textbf{endpoint}, and since the graph only
travels in one \textbf{direction} from it, that endpoint is always an extreme value. $y=\sqrt{x}$
climbs away from $(0,0)$, so the endpoint is the lowest point --- an absolute minimum. $p$ falls away
from $(0,2)$, so its endpoint is the highest point --- an absolute maximum. $q$ has \emph{no}
endpoint at all and runs to $-\infty$ on the left and $+\infty$ on the right, so it has no extreme
value of either kind.}
\end{tcolorbox}

\begin{tcolorbox}[colback=white, colframe=black!40, title=\textbf{Tier E --- Extension},
    fonttitle=\bfseries, arc=1.5mm, left=3mm, right=3mm, top=2mm, bottom=2mm, breakable]
\textbf{Part 1 --- No graph allowed.} Find each domain from the equation alone.
\begin{center}
\renewcommand{\arraystretch}{1.6}
\begin{tabularx}{\linewidth}{@{}l X X X@{}}
\hline
\textbf{Function} & \textbf{Index: even or odd?} & \textbf{Condition on the radicand} & \textbf{Domain (interval notation)} \\
\hline
$\sqrt{x-7}$ & \ans{even} & \ans{$x-7\ge 0$} & \ans{$[7,\infty)$} \\
$\sqrt[3]{x+5}$ & \ans{odd} & \ans{none needed} & \ans{$(-\infty,\infty)$} \\
$\sqrt{12-3x}$ & \ans{even} & \ans{$12-3x\ge 0$} & \ans{$(-\infty,4]$} \\
$\sqrt{2x+1}$ & \ans{even} & \ans{$2x+1\ge 0$} & \ans{$[-\tfrac{1}{2},\infty)$} \\
\hline
\end{tabularx}
\end{center}
Exactly one of these four graphs opens to the \textbf{left}. Which one, and what in the equation
tells you so \emph{before} you graph anything? \ansline{$\sqrt{12-3x}$. Its radicand has a
\emph{negative} coefficient on $x$, so the radicand only stays non-negative as $x$ gets
\emph{smaller}; the domain $(-\infty,4]$ is a ray heading left, so the graph runs left from its
endpoint at $(4,0)$.}

\vspace{5pt}
\textbf{Part 2 --- A radical model, and the family it came from.} Ignoring air resistance, an object
dropped from a height of $d$ feet hits the ground after
\[
  t(d) = \frac{\sqrt{d}}{4} \ \text{ seconds.}
\]
\begin{center}
\renewcommand{\arraystretch}{1.25}
\begin{tabular}{|c|c|c|c|c|c|}
\hline
$d$ (feet) & $0$ & $16$ & $64$ & $144$ & $256$ \\ \hline
$t(d)$ (seconds) & $0$ & $1$ & $2$ & $3$ & $4$ \\ \hline
\end{tabular}
\end{center}
\begin{enumerate}[label=(\alph*), leftmargin=1.9em, itemsep=6pt]
  \item What is the domain of $t$ \emph{algebraically} (from the radicand)? \ans{$[0,\infty)$} \\
        Does the context agree, or does it restrict things further? \ansline{It agrees exactly. The
        radicand $d$ must be $\ge 0$, and a drop height cannot be negative either --- this is the
        rare case where the algebra and the context give the same answer.}
  \item The graph of $t$ has an endpoint at $(0,0)$. What does that point \emph{mean} for a falling
        object? \ansline{Dropped from a height of $0$ feet, the object is already on the ground, so
        it takes $0$ seconds to fall. The endpoint is the ``no drop at all'' case.}
  \item From the table: going from $d=64$ to $d=256$ multiplies the height by \ans{4}, but
        multiplies the falling time by only \ans{2}. Use the shape of a square root graph to
        explain why the time grows so much more slowly than the height. \ansline{A square root graph
        flattens as it climbs, so equal jumps in output require ever-larger jumps in input. Because
        square-rooting undoes squaring, you must multiply the height by $4$ to double the time --- by
        $9$ to triple it, and so on.}
  \item Solve $t=\dfrac{\sqrt{d}}{4}$ for $d$ in terms of $t$. \quad $d=$ \ans{$16t^2$} \\
        Which function family is that? \ans{quadratic} \quad What does this tell you about the
        relationship between $t(d)$ and the function you just wrote? \ansline{They are
        \textbf{inverses} --- each undoes the other. One takes a height and returns a time; the other
        takes a time and returns a height. Their graphs are reflections over $y=x$, and their domains
        and ranges are swapped.}
  \item Your answer to (d) is a parabola, but only \emph{half} of it describes falling objects.
        Which half, and why? Connect this to the reason $y=x^2$ needed a \textbf{restricted domain}
        in the notes. \ansline{Only the half with $t\ge 0$, since negative time is meaningless here.
        That is the same restriction the notes put on $y=x^2$: the full parabola sends two inputs
        ($t$ and $-t$) to the same height, so it could not be turned back into a single function
        until one half was thrown away. Restricting to $t\ge 0$ is exactly what makes $16t^2$ and
        $\frac{\sqrt{d}}{4}$ a true inverse pair.}
\end{enumerate}
\end{tcolorbox}

\begin{teachernote}
\textbf{Tier R} exists for one sentence: item~4. Students who have only ever met functions defined
everywhere assume every graph has a $y$-intercept. Watching $g(0)=\sqrt{-1}$ fail is the cheapest way
to make ``restricted domain'' mean something. Item~6 is the second: accept only ``the graph stops,''
never a made-up left arm.

\textbf{Tier A} is the assessed core. The two graphs are chosen so that \emph{every} row of the table
differs --- there is no answer a student can carry across by habit. The likeliest errors are giving
$p$ the range $[2,\infty)$ (reading the endpoint as a minimum out of reflex) and giving $q$ an
endpoint at $(0,-1)$ (mistaking the labeled $y$-intercept for the start of the graph). Both are worth
a whole-class stop.

\textbf{Tier E} carries the unit's forward links. Part~1 row~3 is the sign-flip from Warm-Up item~2(c)
and the source of the ``opens left'' answer; expect $x\ge 4$ from at least one group. Part~2 is the
Lesson~5.7 seed --- students derive an inverse pair from a real context and rediscover the domain
restriction on their own, which is a far stronger memory than being told. If a group finishes early,
ask them what the \emph{third} member of the pattern is: what would you have to restrict for
$y=x^4$, and why does $y=x^5$ need nothing?
\end{teachernote}

\end{document}
